\chapter{Lý sự cùn sát lớp thật}
\markboth{Lý sự cùn sát lớp thật}{Lý sự cùn sát lớp thật}

\section{Em không cãi, em chỉ không phục}
\begin{truyen}{Em không cãi, em chỉ không phục}{Classroom Talk}
\chuhoa{M}{ột bạn bị nhắc thì bật lại ngay: ``Em không cãi đâu. Em chỉ không phục cách thầy nói thôi.''}

Thầy giữ giọng bình tĩnh: ``Không phục thì em có quyền nói. Nhưng nói để hiểu ra khác với nói để giành phần thắng.''

Bạn ấy vẫn cứng: ``Thì em đang nói đó.''

Thầy đáp: ``Ừ. Và cả lớp đang nghe xem em muốn làm rõ vấn đề hay chỉ muốn giữ sĩ diện.''

\textit{(Disagreement can be healthy, but the tone often reveals whether someone wants understanding or victory.)}
\end{truyen}

\begin{baihoc}
Không phải cứ to tiếng là có lý. Nhiều khi chỉ là cái tôi đang bị dồn góc.
\textit{(A louder tone does not prove stronger logic. Sometimes it only reveals a cornered ego.)}
\end{baihoc}

\begin{ghinhoanh}
\textbf{Short English to remember:}

Disagreement is allowed.

Ego is often louder than reason.
\end{ghinhoanh}

\begin{tuhocnhanh}
1. Khi em phản ứng mạnh, em đang bảo vệ lẽ phải hay sĩ diện? \textit{(When you react strongly, are you defending truth or pride?)}

2. Em có biết nói khác ý mà vẫn giữ tôn trọng không? \textit{(Can you disagree while staying respectful?)}
\end{tuhocnhanh}
\ngancach

\section{Em ngủ nhưng đầu óc em vẫn online}
\begin{truyen}{Em ngủ nhưng đầu óc em vẫn online}{Classroom Talk}
\chuhoa{C}{ả lớp cười khi một bạn đang gục đầu bị gọi dậy liền đáp: ``Mắt em nghỉ thôi, chứ đầu óc em vẫn online.''}

Thầy hỏi: ``Online vậy thầy vừa giảng tới đâu?''

Bạn ấy mở mắt một hồi rồi nói: ``Dạ... tới phần em cần tỉnh rồi.''

Thầy cười: ``Đúng. Phần đó bắt đầu từ ba phút trước.''

\textit{(Humorous excuses often collapse under one simple factual question.)}
\end{truyen}

\begin{baihoc}
Nói vui có thể cứu một pha quê. Nhưng nó không cứu được lỗ hổng vừa xảy ra.
\textit{(A joke may save a moment of embarrassment, but it does not repair the loss that just happened.)}
\end{baihoc}

\begin{ghinhoanh}
\textbf{Short English to remember:}

Humor can lighten the moment.

It cannot replace attention.
\end{ghinhoanh}

\begin{tuhocnhanh}
1. Em đang thiếu ngủ thật hay thiếu kỷ luật ngủ nghỉ? \textit{(Are you truly sleep-deprived, or lacking sleep discipline?)}

2. Em cần sửa gì để vào lớp tỉnh hơn? \textit{(What needs to change so you enter class more awake?)}
\end{tuhocnhanh}
\ngancach

\section{Em không quay bài, em chỉ nhìn cho đỡ hoảng}
\begin{truyen}{Em không quay bài, em chỉ nhìn cho đỡ hoảng}{Classroom Talk}
\chuhoa{M}{ột bạn bị bắt gặp nhìn ngang liền chống chế: ``Em không quay bài. Em chỉ liếc qua cho đỡ panic thôi.''}

Thầy hỏi: ``Vậy nếu em chỉ lấy tiền người ta cho đỡ hoảng, có gọi là lấy không?''

Cả lớp ồ lên. Bạn kia cười gượng: ``Thầy ví mạnh tay quá.''

Thầy đáp: ``Vì em đang làm nhẹ một chuyện không nhẹ.''

\textit{(Students often soften the language around cheating, but changing the label does not change the act.)}
\end{truyen}

\begin{baihoc}
Đổi tên một hành vi không làm hành vi đó sạch hơn.
\textit{(Renaming an action does not make it cleaner.)}
\end{baihoc}

\begin{ghinhoanh}
\textbf{Short English to remember:}

A softer label does not change the act.

Cheating remains cheating.
\end{ghinhoanh}

\begin{tuhocnhanh}
1. Em có hay làm nhẹ lỗi bằng cách đổi cách gọi không? \textit{(Do you soften your mistakes by renaming them?)}

2. Khi hoảng, em có cách nào khác ngoài nhìn bài người khác? \textit{(When you panic, what can you do besides looking at someone else’s paper?)}
\end{tuhocnhanh}
\ngancach

\section{Em không hỗn, em nói thẳng thôi}
\begin{truyen}{Em không hỗn, em nói thẳng thôi}{Classroom Talk}
\chuhoa{C}{ó bạn phản ứng rất gắt: ``Em không hỗn. Em chỉ nói thẳng. Người lớn cứ thích vòng vo nên nghe nói thẳng mới khó chịu.''}

Thầy đáp: ``Nói thẳng là nói đúng và rõ. Còn đâm người khác bằng lời rồi gọi đó là thẳng tính thì dễ quá.''

Bạn ấy chậm lại: ``Vậy sao phân biệt?''

Thầy nói: ``Nói xong mà vấn đề sáng hơn, đó là thẳng. Nói xong chỉ làm người khác nghẹn lại, thường là thiếu tôn trọng.''

\textit{(Directness is not cruelty. If a sentence clarifies the issue, it helps. If it only wounds, it fails.)}
\end{truyen}

\begin{baihoc}
Thẳng tính không phải giấy phép để thiếu lịch sự.
\textit{(Being direct is not a license for disrespect.)}
\end{baihoc}

\begin{ghinhoanh}
\textbf{Short English to remember:}

Direct is not the same as rude.

Clarity should not crush respect.
\end{ghinhoanh}

\begin{tuhocnhanh}
1. Em có đang dùng chữ ``thẳng'' để che một cách nói làm đau người khác không? \textit{(Do you use the word direct to hide hurtful speech?)}

2. Em có thể nói cùng ý đó theo cách nào tốt hơn? \textit{(How could you say the same thing in a better way?)}
\end{tuhocnhanh}
\ngancach